% ============================================================
% ch10.tex —— 第 10 章 环与域：什么时候才能做除法
% ============================================================
\chapter{环与域：什么时候才能做除法}

第 5 章留下过一张支票：在钟表世界 $Z_{12}$ 里，$2$ 没有倒数，"除以 $2$"是非法操作。这一章正式兑付——顺带回答三个更深的问题：除法全面通行的世界长什么样？素数在其中扮演什么角色？以及，\textbf{分数这种东西，究竟是怎么被人类制造出来的？}（答案：一条四步流水线，全程只用整数。）

\section{两种运算并肩干活的世界}

先看清楚我们脚下已有哪些世界。整数世界的好，好在加法与乘法\textbf{并肩工作}：各自守规矩（加法是群、乘法封闭且结合有单位元），再用\textbf{分配律} $a\times(b+c) = a\times b + a\times c$ 把两座桥焊起来。分配律不是小事——它保证了"提公因式"合法（$ab + ac = a(b+c)$，你做了几百道因式分解题的许可证）。把这种"双运算结构"抽象出来：

\begin{dingyi}[环]
一个集合配两种运算 $+$ 和 $\times$，若满足：
\begin{enumerate}
  \item 加法四项体检全过（构成群，且交换）；
  \item 乘法封闭、满足结合律、有单位元 $1$；
  \item 分配律连接两运算：$a\times(b+c) = ab + ac$，$(b+c)\times a = ba + ca$；
\end{enumerate}
则称为一个\textbf{环}（ring）。
\end{dingyi}

点名报到：\textbf{整数} $\mathbb{Z}$ 是环（原型机）；\textbf{钟表世界} $Z_n$ 是环（加法绕圈、乘法压扁——分配律验证：$(a+b)$ 先压扁再乘 $c$，与先各自乘再压扁，结果相同，可随手抽三个数试）；\textbf{多项式}是环（下一章主角）。环的世界里，加减乘畅通无阻——唯独\textbf{除法不保证}。整数环里 $1\div 2$ 就卡住了。

\section{除法的通行证审查}

在环里做除法 $a \div b$，本质是解方程 $b\times x = a$。而它的总闸门是更基础的问题：\textbf{$b\times x = 1$ 有没有解}——即 $b$ 有没有\Keyword{倒数}（正式名：逆元；有倒数的元素叫\Keyword{单位}，unit）。有倒数，除法就能化成乘倒数，全线畅通；没倒数，$b$ 就是"除法路障"。逐个城市审查：

\paragraph{$Z_{12}$：部分通行。}把 $12$ 位居民的乘法倒查表列出来（找"$xy \equiv 1$"的搭档）：
\[
\begin{array}{llll}
1\times1 = 1 & 5\times5 = 25\equiv 1 & 7\times7 = 49\equiv 1 & 11\times 11 = 121 \equiv 1
\end{array}
\]
\textbf{有倒数的名单：$\{1, 5, 7, 11\}$}——恰好是与 $12$ 互素的四个数（$\gcd(1,12)=\gcd(5,12)=\gcd(7,12)=\gcd(11,12)=1$）。没倒数的 $2,3,4,6,8,9,10$，每个都与 $12$ 有大于 $1$ 的公因数。规律若隐若现。

\paragraph{$Z_4$：惨案现场。}列全乘法表（$0,1,2,3$）看对角线之外：$2\times2 = 4 \equiv 0$——\textbf{两个非零数相乘，居然归零}！第 1 章烧脑时刻埋的雷在这里引爆。这种元素（非零、但能和另一个非零数乘出零）叫\Keyword{零因子}。零因子是除法世界的头号通缉犯，一经定罪，立即注销倒数资格——证明只需三行（反证法，第 4 章的战术）：设 $b$ 是零因子，$b\times c = 0$（$c \ne 0$），假设 $b$ 有倒数 $x$，则
\[
c = 1\times c = (x\times b)\times c = x\times(b\times c) = x\times 0 = 0 ,
\]
与 $c\ne0$ 矛盾。$b$ 无倒数，证毕。（这里用到"$x\times0=0$"，序章习题你用分配律证过。）

\paragraph{$Z_5$：全面通行。}五位居民的倒数名单：$1^{-1} = 1$，$2^{-1} = 3$（$2\times3 = 6\equiv1$），$3^{-1} = 2$，$4^{-1} = 4$（$16\equiv1$）——\textbf{非零居民人人有倒数}。零因子？想找都找不到（两个 $1$--$4$ 的数相乘，模 $5$ 永远不是 $0$）。

三个城市的命运对照，指向一条定理。$5$ 是素数，$4$、$12$ 是合数——是巧合吗？

\begin{dingli}[素数 $=$ 除法通行证]
$Z_n$ 中每个非零数都有倒数 $\iff$ $n$ 是素数。
\end{dingli}

\begin{proof}
$\Leftarrow$（素数 $\Rightarrow$ 全员有倒数）：设 $p$ 素数，$a \in \{1, \dots, p-1\}$。则 $\gcd(a, p) = 1$（$p$ 的因数只有 $1$ 和 $p$，$p\nmid a$，公共因数只剩 $1$）。第 2 章裴蜀等式：存在整数 $x, y$ 使 $ax + py = 1$。两边取模 $p$（第 1 章替换原理，$py$ 那项压成 $0$）：
\[
ax \equiv 1 \pmod p ,
\]
$x\bmod p$ 就是 $a$ 的倒数。\textbf{两千年前铺地砖的辗转相除，在这里交出了每一个倒数}——不光证明它存在，还给出计算法（裴蜀系数）。

$\Rightarrow$（合数 $\Rightarrow$ 有人无倒数）：设 $n = ab$（$1 < a, b < n$，合数定义）。在 $Z_n$ 里 $a\times b = n \equiv 0$：$a$ 是零因子。零因子无倒数（上面三行反证法），除法破功。证毕。
\end{proof}

素数在这里完成了第二次登基：第 3 章它是"乘法的原子"（唯一分解），本章它是"\textbf{除法的通行证}"（$Z_p$ 无零因子、人人可倒）。同一个对象，两种视角，互相印证——这正是抽象的复利。同时回收第 6 章的伏笔："模素数的世界里约去合法"，因为约去 $a$ 就是乘 $a^{-1}$，而 $Z_p$ 里逆元人人有。

\begin{dingyi}[域]
一个环，若它的\textbf{非零元素全体}在乘法下构成群（即人人有倒数、无零因子），就升级称为\textbf{域}（field）。
\end{dingyi}

\begin{cidian}
土名对照：并肩干活的世界 $=$ \textbf{环}（ring）；除法全面通行的环 $=$ \textbf{域}（field）；假零制造者 $=$ \textbf{零因子}（zero divisor）；有倒数的数 $=$ \textbf{单位}（unit）。你小学就认识两个域：分数世界 $\mathbb{Q}$、实数世界 $\mathbb{R}$；现在又认识一整族袖珍域 $Z_2, Z_3, Z_5, Z_7, Z_{11}, \dots$——注意 $Z_4, Z_6, Z_8, Z_9, Z_{10}, Z_{12}$ 都\textbf{不是}域。
\end{cidian}

\section{流水线现场：分数是怎样炼成的}

域这么好用，一个自然的野心升起：\textbf{能不能把任何"没倒数"的环，加工成域？}以整数环为原料，现场直播。痛点明确：$2\div3$ 在整数里没结果——世界上没有哪个整数乘 $3$ 等于 $2$。\textbf{发明许可证}掏出来：既然这东西不存在，就造一个。

\textbf{第一步：宣告新数的形态}。新数是一个"\textbf{整数对}" $(a, b)$（$b \ne 0$），直觉含义是"$a$ 份 $b$"，即 $a\div b$ 的化身。例：$(6,2)$ 意为 $6\div2$，$(2,3)$ 意为 $2\div3$。

\textbf{第二步：规定"哪些对算同一个数"}。直觉说 $(6,2)$ 和 $(3,1)$ 是同一个数（都等于 $3$）。但"直觉"不能当法律，法律必须用整数写：
\[
(a, b) \sim (c, d) \;\iff\; a\times d = c\times b .
\]
检验：$(6,2)\sim(3,1)$？$6\times1 = 6$，$3\times2 = 6$，相等，\checkmark 同一个数。$(2, 4)\sim(3,6)$？$2\times6 = 12 = 3\times4$ \checkmark（都是"一半"）。而 $(1,2) \sim (1,3)$？$1\times3 = 3 \ne 2 = 1\times2$，不同数 \checkmark。这条规定就是初中"交叉相乘验比例"（$\tfrac ab = \tfrac cd \iff ad = bc$）的机器版。

\textbf{第三步：定义运算}（设计的唯一目标是：让新世界继承旧世界的一切好规矩）：
\[
(a,b) + (c,d) := (ad+cb,\; bd), \qquad (a,b)\times(c,d) := (ac,\; bd).
\]
以 $\tfrac12 + \tfrac13$ 试机：$(1,2)+(1,3) = (1\times3 + 1\times2,\ 2\times3) = (5, 6)$——正是通分加法 $\tfrac12+\tfrac13 = \tfrac56$。乘法 $\tfrac23\times\tfrac34$：$(2,3)\times(3,4) = (6, 12) \sim (1,2)$——分子乘分子、分母乘分母，再化简。\checkmark 机器正常。

\textbf{第四步：出厂质检}（四项体检 + 两条"良定义"）。
\begin{itemize}
\item 加法单位元 $(0,1)$：$(a,b) + (0,1) = (a\times1 + 0\times b,\ b\times1) = (a, b)$ \checkmark；
\item 加法逆元 $(-a, b)$；乘法单位元 $(1,1)$；
\item \textbf{倒数梦想成真}：非零数 $(a,b)$（$a\ne0$）的乘法逆元是 $(b,a)$——$(a,b)\times(b,a) = (ab, ab) \sim (1,1)$ \checkmark。整数世界里 $2$ 找不到搭档的千年之憾，在新世界一行解决；
\item 分配律：机械展开可验（两边都是三元组，逐项对上）；
\item \textbf{良定义检查}（流水线特有的质检）：把 $(6,2)$ 换成等价的 $(3,1)$ 参与运算，答案必须不变。抽验：$(6,2)+(1,3)$ 按 $(6,2)$ 算得 $(6\times3+1\times2,\ 6) = (20, 6) \sim (10, 3)$；按 $(3,1)$ 算得 $(3\times3+1,\ 3) = (10,3)$。\checkmark 同一个答案。（"等价的写法算出等价的答案"必须对\textbf{所有}运算成立，本书抽验了加法，乘法同理。）
\end{itemize}
质检通过，\textbf{一个崭新的域下线}。它就是分数世界 $\mathbb{Q}$（有理数域）。你写了十几年的"$\tfrac23$"，从今天起有了出生证明：\textbf{它是"整数对 $(2,3)$"经四步流水线锻造成的域成员}。"分数自带两套写法（$\tfrac24$ 与 $\tfrac12$）"也不再是讨厌的特例，而是"等价类"的日常——一个数，多名身份证。

这台流水线（数学家叫它\textbf{局部化}或\textbf{分式域构造}）还能开第二次：把多项式环（下一章）喂进去，产出"多项式分数"$\tfrac{x+1}{x^2+3}$——微积分里分式函数的正式户口。第 17、18 章的积分表里住着大量这种居民。

\section{最小的域：$Z_2$ 与信息时代}

域不一定要大。$Z_2 = \{0, 1\}$ 只有两位居民，加法 $1+1 = 0$，乘法照旧。它虽小，五脏俱全（$1$ 的倒数是 $1$）。而它统治着你的生活：\textbf{计算机的每个比特就是 $Z_2$ 的居民}。电压高/低、磁极 N/S、开关开/关——全部是 $0/1$；"异或"运算就是 $Z_2$ 加法（相同得 $0$，不同得 $1$）、"与"运算就是 $Z_2$ 乘法。纠错码（二维码破了角也能扫）、CDMA 通信、AES 加密，底层全是 $Z_2$ 上的代数。\textbf{最小的域，扛着最大的产业}。

\begin{dongshou}
\begin{itemize}
  \yisi 列出 $Z_7$ 的完整倒数表（$1^{-1}=1$，$2^{-1}=4$，$3^{-1}=5$，补 $4,5,6$；注意表是对称的：$x^{-1} = y \iff y^{-1} = x$）。
  \yier 找出 $Z_6$ 的全部零因子（$2, 3, 4$），并验证零因子名单 $=$ "与 $6$ 不互素的非零数"名单。再问：$Z_6$ 里 $3\times x \equiv 0$ 有几个解？（$x = 0, 2, 4$——三个，因为 $\gcd(3,6) = 3$ 把解撑成 $3$ 个。第 5 章"解不止一个"的怪现象，全部病灶都在这。）
  \yier 用流水线记号演算 $(2,3) + (1,6)$ 与 $(2,3)\times(3,4)$，并与 $\tfrac23+\tfrac16 = \tfrac56$、$\tfrac23\times\tfrac34 = \tfrac12$ 对照。
  \yisi 在 $Z_2$ 里解 $x^2 + x = 0$。（代入 $x=0$：$0$ \checkmark；$x=1$：$1+1 = 0$ \checkmark——\textbf{两个解都是根}。对比实数世界 $x^2+x=0$ 的根 $0, -1$：不同的域，不同的居民，一样的"完全平方分解"$(x)(x+1)$——在 $Z_2$ 里 $x+1 = x-1$，因为 $\pm1$ 不分家。）
\end{itemize}
\end{dongshou}

\begin{shaonao}
\textbf{费马小定理的"除法版"证明}。目标：$p$ 素数时 $\binom{p}{k} = \dfrac{p!}{k!(p-k)!}$（组合数，$0<k<p$）\textbf{必被 $p$ 整除}。思路：$p! = 1\times2\times\cdots\times p$ 含因子 $p$，而分母 $k!(p-k)!$ 的成员全部小于 $p$、与 $p$ 互素——所以分母的每个成员都在 $Z_p$ 里有倒数，除法在 $Z_p$ 里合法。于是在 $Z_p$ 里 $\binom pk = 0$，即 $p \mid \binom pk$。用 $p = 5$ 验证：$\binom51 = 5$，$\binom52 = 10$，$\binom53 = 10$，$\binom54 = 5$，全是 $5$ 的倍数 \checkmark。这个"系数全被 $p$ 整除"的事实，正是 $(a+b)^p \equiv a^p + b^p \pmod p$ 的原因（中间系数全部蒸发）——两边再迭代，就是"新鲜版"费马小定理的第三种证明（第一种：第 6 章幂循环；第二种：第 9 章拉格朗日；\textbf{一鱼三吃}）。
\end{shaonao}

\begin{chengshi}
环的教科书定义还要求加法交换、部分教材不要求乘法单位元（两种流派），我们采用"够用版"。分数流水线里"$\sim$ 是等价关系（自反、对称、传递）"与"良定义对一切运算成立"的完整验证是纯机械的交叉相乘练习（传递性最费笔：$ad=cb,\ ce=fd \Rightarrow ae = fb$，两边同乘适当因子），本书演示了操作、略去了地毯下的例行公事。"$Z_p$ 是域"的乘法逆元存在性依赖费马小定理\textbf{或}裴蜀（我们选了裴蜀，因为它还给计算法）；域的完整分类（有限域不只 $Z_p$，还有 $p^k$ 元的 GF$(p^k)$——第 32 章的椭圆曲线密码就工作在 $p^k$ 元域上）超出本书，但"有限域必有无穷多"这条知道一下很提气。
\end{chengshi}

环里最能干的员工，下一章登场：多项式。你会发现它和整数像到令人怀疑人生——所有整数定理，翻译一遍全部免费重装。
